%%% FIELD proof-bandwidth-feasibility
By \(\text{ass:polynomial-bandwidth}\), write \(h_n=n^{-a_h}\) with \(a_h>0\), and by the definition of \(L_n\), \(L_n=a_h\log n\).  Since \(s\geq 8\), all deterministic powers below have positive smoothness exponents.

For the first displayed requirement, direct substitution gives
\[
 \sqrt{\frac{n}{L_n}}h_n^{s-2}
 =\frac{1}{\sqrt{a_h\log n}}\,n^{1/2-a_h(s-2)}.
\]
If \(a_h>1/[2(s-2)]\), the power of \(n\) is negative, and if equality holds, the expression is \((a_h\log n)^{-1/2}\).  Both cases converge to zero.  If \(a_h<1/[2(s-2)]\), the positive power of \(n\) dominates the logarithmic denominator, so the expression diverges.  Thus the first requirement holds exactly when
\[
 a_h\geq \frac{1}{2(s-2)}.
\]

For the second requirement, its stochastic term is
\[
 \sqrt{\frac{\log n}{n h_n^{13/2}}}
 = (\log n)^{1/2}n^{-1/2+(13/4)a_h}.
\]
It converges to zero exactly when \(-1/2+(13/4)a_h<0\), namely when \(a_h<2/13\); at equality it equals \((\log n)^{1/2}\), and above equality it diverges.  Its deterministic term satisfies
\[
 h_n^{s-9/4}=n^{-a_h(s-9/4)}\longrightarrow0
\]
because \(a_h>0\) and \(s-9/4>0\).  Since both summands are nonnegative, the second requirement therefore holds exactly when \(a_h<2/13\).

For the third requirement, put
\[
 A_n=\sqrt{\frac{\log n}{n h_n^{9/2}}},
 \qquad
 B_n=h_n^{s-5/4}.
\]
The three nonnegative terms in \(\sqrt{n/L_n}(A_n+B_n)^2\) are
\[
 \sqrt{\frac{n}{L_n}}A_n^2
 =\frac{(\log n)^{1/2}}{\sqrt{a_h}}
   n^{-1/2+(9/2)a_h},
\]
\[
 2\sqrt{\frac{n}{L_n}}A_nB_n
 =\frac{2}{\sqrt{a_h}}n^{-a_h(s-7/2)},
\]
and
\[
 \sqrt{\frac{n}{L_n}}B_n^2
 =\frac{1}{\sqrt{a_h\log n}}
   n^{1/2-a_h(2s-5/2)}.
\]
The first of these converges to zero exactly when \(a_h<1/9\); at \(a_h=1/9\) it diverges as \((\log n)^{1/2}\).  The mixed term converges to zero for every \(a_h>0\), since \(s-7/2>0\).  The last term converges to zero exactly when
\[
 a_h\geq \frac{1}{4s-5};
\]
equality is admitted because then only the factor \((a_h\log n)^{-1/2}\) remains.  Nonnegativity of the three terms also proves necessity.  Hence the third requirement holds exactly when \(a_h<1/9\) and \(a_h\geq1/(4s-5)\).

It remains to intersect the restrictions.  For \(s\geq8\),
\[
 \frac{1}{2(s-2)}>\frac{1}{4s-5},
 \qquad
 \frac{1}{9}<\frac{2}{13},
 \qquad
 \frac{1}{2(s-2)}<\frac{1}{9}.
\]
Consequently the lower restriction from the first requirement implies the deterministic-square restriction, while the upper restriction from the stochastic square implies the pilot restriction.  The conjunction of all three displayed scale requirements is therefore exactly
\[
 \frac{1}{2(s-2)}\leq a_h<\frac{1}{9}.
\]
The equality at the lower endpoint is valid because the order-minus-two bias is then divided by \(\sqrt{L_n}\), whereas the upper endpoint is excluded because the normalized stochastic square retains the divergent factor \((\log n)^{1/2}\).  This proves the claim.
